\subsection{C.4 continued: construction without a previously realized group}

The central-extension parameter in scans 47--49 is a torsor parameter:
there is no preferred origin until a compatible factor set is chosen.
We now compute that factor set directly from the given crossed modules.
No group realizing the prescribed centre and derived group is assumed.
The result in this subsection is about ordinary groups, as is
Theorem~\ref{C:central_twists}.  Constant group schemes over an
algebraically closed field, or over a perfect field, satisfy the same
criterion, in every characteristic.  A statement about arbitrary formal
or algebraic group schemes additionally requires their sheaf-theoretic
realization and representability; the ordinary-group assertion does not
silently supply that additional assertion.

\begin{theorem}[Proved: an explicit realization criterion without a baseline]
\label{C:baseline-free-realization}
Suppose we are given groups and exact sequences
\[
 1\longrightarrow C\longrightarrow N\xrightarrow{d}L
       \xrightarrow{q_L}Q\longrightarrow1,\qquad
 1\longrightarrow C\longrightarrow Z\xrightarrow{e}M
       \xrightarrow{q_M}Q\longrightarrow1.
\]
The injections of $C$ are part of the data.  Assume that $Z$ and $M$
are abelian, that $d:N\to L$ is a crossed module for an action
$\theta:L\to\operatorname{Aut}(N)$, and that $L$ acts trivially on $C$.
Thus
\[
 d(\theta_l(n))=l d(n)l^{-1},\qquad
 \theta_{d(n)}(n')=nn'n^{-1}.
\]
In particular, $C$ and $Q$ are abelian.  Put
\[
 H=(N\times Z)/\{(c,c^{-1}):c\in C\},\qquad K=L\times_Q M.
\]
We regard $N,Z$ as subgroups of $H$, with intersection $C$.
Define $\partial:H\to K$ by
$\partial(nz)=(d(n),e(z))$.  Let $(l,m)\in K$ act on $H$ by
$nz\mapsto\theta_l(n)z$.  This is a crossed module, with kernel $C$
and cokernel $Q$.

Here is a complete test, using only these data, for all marked groups
$G$ with $N\triangleleft G$, $Z\subset Z(G)$, $G/Z=L$, $G/N=M$,
and the prescribed maps and actions, followed by the two equalities
$Z=Z(G)$ and $N=G'$.
Choose a normalized set section $s:Q\to K$, and choose normalized
$f:Q\times Q\to H$ such that
\[
 \partial f(x,y)=s(x)s(y)s(xy)^{-1}.
\]
Write $a_x$ for the automorphism of $H$ induced by $s(x)$, and compute
\[
 o(x,y,z)=a_x(f(y,z))f(x,yz)f(xy,z)^{-1}f(x,y)^{-1}\in C.
                                                        \tag{C.4.3}
\]
This is a normalized $3$-cocycle for the trivial action of $Q$ on $C$.
There is a marked realization if and only if the explicit equation
\[
 (\delta t)(x,y,z)
 =t(y,z)t(x,yz)t(xy,z)^{-1}t(x,y)^{-1}=o(x,y,z)^{-1}
                                                        \tag{C.4.4}
\]
has a normalized $C$-valued $2$-cochain solution.  Equivalently,
$[o]=0$ in $H^3(Q,C)$.

If it does, fix one solution $t$ and put $F=tf$.  For each normalized
$2$-cocycle $c:Q^2\to C$, put $F_c=cF$ and define on $H\times Q$
\[
 (h,x)(h',y)=\bigl(h\,a_x(h')\,F_c(x,y),xy\bigr).
                                                        \tag{C.4.5}
\]
This is a marked realization $G_c$.  Every marked realization occurs,
and two such cocycles give the same marked isomorphism class exactly
when their quotient is a $2$-coboundary.  Hence the marked realizations,
when nonempty, form a torsor under $H^2(Q,C)$.

The following tests determine exactly the admissible central extensions
$E_c$ of $Q$ by $C$.  They contain no unknown group $G$.
Put
\[
 \begin{split}
 W&=Z(K)\cap\ker\bigl(K\longrightarrow\operatorname{Aut}(H)\bigr),\qquad
 \zeta=W/\partial Z,\\
 J&=\bigl\langle [n,n'],\ a_x(n)n^{-1}:n,n'\in N,\ x\in Q\bigr\rangle
       \subset N,\qquad \eta=N/J.
 \end{split}                                             \tag{C.4.6}
\]
Here $J$ is equivalently the smallest normal, $a_x$-stable subgroup
containing the displayed elements.  Write
$\alpha:\zeta\to Q$ and $\beta:C\to\eta$ for the induced maps.
For $w\in W$, choose $x\in Q$ and $h\in H$ with $w=\partial h\,s(x)$.
Then the formulas
\[
 \begin{split}
 \psi_F(w\bmod\partial Z)(y)
   &=hF(x,y)\bigl(a_y(h)F(y,x)\bigr)^{-1}\in C,\\
 \lambda_F(x\wedge y)&=F(x,y)F(y,x)^{-1}\bmod J\in\eta
 \end{split}                                             \tag{C.4.7}
\]
define homomorphisms
$\psi_F:\zeta\to\operatorname{Hom}(Q,C)$ and
$\lambda_F:\Lambda^2_{\mathbb Z}Q\to\eta$.
For $b_c(x,y)=c(x,y)c(y,x)^{-1}$, the answer is exactly
\[
 \boxed{\begin{aligned}
 Z(G_c)=Z&\iff
   \psi_F+\bigl[z\mapsto b_c(\alpha(z),-)\bigr]\text{ is injective},\\
 G_c'=N&\iff \lambda_F+\beta b_c\text{ is surjective}.
 \end{aligned}}                                         \tag{C.4.8}
\]
The sums in this display use additive notation for the abelian targets.
One may equivalently test the centre without $\zeta$: every $(h,x)$
satisfying
\[
 \operatorname{Int}(h)a_x=1\text{ on }H,\qquad
 hF_c(x,y)=a_y(h)F_c(y,x)\quad(y\in Q)
                                                        \tag{C.4.9}
\]
must have $x=1$ and $h\in Z$.
\end{theorem}
\begin{proof}
The crossed-module identities imply $C\subset Z(N)$, so the central
product $H$ is defined.  Its maps from $N$ and $Z$ are injective and
have intersection $C$: for example $(n,1)=(1,z)$ in the quotient
holds exactly when $n=z\in C$.  The action on $H$ is well-defined
because $L$ fixes $C$, and the crossed-module identities for $\partial$
follow from those for $d$; the $M$ coordinate is harmless because $M$
is abelian.  If $\partial(nz)=1$, then $n,z\in C$, so its kernel is
the diagonal copy of $C$ in $H$.  Its image is
$\ker(K\to Q)=(N/C)\times(Z/C)$, proving the cokernel assertion.

The elements $a_x$ and $f$ satisfy
\[
                  a_xa_y=\operatorname{Int}(f(x,y))a_{xy}.
                                                        \tag{C.4.10}
\]
Applying $\partial$ to the two products
$a_x(f(y,z))f(x,yz)$ and $f(x,y)f(xy,z)$ gives the same element
$s(x)s(y)s(z)s(xyz)^{-1}$ of $K$.  Their quotient is therefore
$o(x,y,z)\in C$, and it is central in $H$.  The action on it is trivial.
The cocycle identity can now be checked without commutative assumptions
on $N$: expand the parenthesized product of four symbols
$s(x),s(y),s(z),s(w)$ using (C.4.10).  The route
$((xy)z)w\to(x(yz))w\to x((yz)w)\to x(y(zw))$
changes the factor in $H$ by
\[
                    o(x,y,z)o(x,yz,w)o(y,z,w),
\]
whereas the route
$((xy)z)w\to(xy)(zw)\to x(y(zw))$
changes it by
\[
                    o(xy,z,w)o(x,y,zw).
\]
Both routes expand to the same product of the $a_x$ and $f$; (C.4.10)
moves each inner automorphism across the next factor, and centrality
of $o$ permits the displayed scalar factors to be collected.  This is
exactly $\delta o=1$.

Replacing $f$ by $tf$, with values of $t$ in $C$, changes $o$ to
$o\delta t$.  If a marked realization exists, lift each $s(x)$ to
$G$ and compute its factor set.  That factor set differs from $f$
by a $C$-valued cochain, and associativity says precisely (C.4.4).
This proves necessity, including the actual comparison map.
Conversely, (C.4.4) makes $F=tf$ satisfy
\[
 a_x(F(y,z))F(x,yz)=F(x,y)F(xy,z).
                                                        \tag{C.4.11}
\]
For (C.4.5), the two products of three pairs have first coordinates
\[
 \begin{split}
 h a_x(h')F_c(x,y)a_{xy}(h'')F_c(xy,z),\\
 h a_x(h')a_xa_y(h'')a_x(F_c(y,z))F_c(x,yz).
 \end{split}
\]
Equations (C.4.10)--(C.4.11) make them equal.  The identity is $(1,1)$;
solving either product equation with $(h,x)$ gives an inverse, and
associativity makes the left and right solutions equal.  Thus this is
a group.  The maps
\[
 H\longrightarrow G_c,\quad h\longmapsto(h,1),\qquad
 G_c\longrightarrow K,\quad (h,x)\longmapsto\partial h\,s(x)
\]
are, respectively, injective and surjective homomorphisms; the second
has kernel $C$ and restricts to $\partial$.  The compositions to $L$
and $M$ have kernels $Z$ and $N$, respectively, giving the prescribed
quotients and action.  Conjugation on $Z$ is trivial by construction.

Every realization is put into this form by the chosen lifts of $s$.
The quotient of two associative factor sets is a $C$-valued $2$-cocycle.
A marked isomorphism fixes $H$ and $K$, so sends each $(1,x)$ to
$(u(x),x)$ with $u(x)\in C$.  Substitution in (C.4.5) says exactly
that the two $2$-cocycles differ by the coboundary of $u$ (or its
inverse according to the direction of the isomorphism).  This proves
both directions of the classification.  Changing the section by
$s'(x)=\partial b(x)s(x)$ replaces the lift of $f$ by
\[
 b(x)a_x(b(y))f(x,y)b(xy)^{-1};
\]
its associator is the same $o$.  Any other lift differs by a central
cochain and hence by a coboundary at the associator level.  Thus the
obstruction class and the torsor of solutions do not depend on choices.
This is the ordinary-group realization obstruction for a crossed module;
we have also opened Brown--Mucuk, Theorem 5.2
\cite{C4-BrownMucuk1994}, whose $M,P,A,Q,\theta$ here are
$H,K,C,Q,\operatorname{id}_Q$.

It remains to impose the two equalities, rather than merely the
inclusions.  An element of $G_c$ centralizing every element has image
in $Z(K)$, and its conjugation on $H$ is trivial.  Conversely an element
above $W$ centralizes $H$, and all of its remaining commutators lie in
$C$.  They therefore give a homomorphism $Q\to C$.  For the element
$(h,x)$ above $w$, its commutator with $(1,y)$ is the first formula
in (C.4.7), with $F_c$ in place of $F$.  This follows by comparing the
two products $(h,x)(1,y)$ and $(1,y)(h,x)$, whose second coordinates
agree because $Q$ is abelian.  Changing the lift of $w$ multiplies it
by $C$, and changing its class modulo $\partial Z$ multiplies its lift
by $Z$; both changes are central.  Hence the formula is well-defined
on $\zeta$.  Central commutator identities give homomorphisms in both
arguments and show
\[
                  Z(G_c)/Z=\ker\psi_{F_c}.
\]
Multiplying $F$ by $c$ multiplies the commutator by
$b_c(x,y)$, giving the first equivalence of (C.4.8), and the same
calculation proves (C.4.9).

The subgroup $[G_c,N]$ is exactly $J$: the group is generated by $N$,
$Z$ and the lifts $(1,x)$; their commutators with $N$ are the generators
listed in (C.4.6), and the normal closure gives both inclusions.
In particular $J$ is independent of $c$, and $\eta=N/J$ is abelian.
Since $G_c/N=M$ is abelian, all commutators lie in $N$.  Modulo $J$,
$N$ is central, so $G_c/J$ has nilpotence class at most two.  The
commutator pairing is alternating and biadditive on $G_c/N=M$.
The image of $Z$ is central, so this pairing factors through
$M/e(Z)=Q$ in each argument.  Comparing the products of $(1,x)$ and
$(1,y)$ gives exactly the second formula in (C.4.7).  Its image is
$G_c'/J$, so it is onto $\eta$ exactly when $G_c'=N$.  Twisting by
$c$ adds $\beta b_c$, proving the second equivalence.  This also proves
that the formulas exhaust all central extensions satisfying both tests.
\end{proof}

\begin{proposition}[Proved: the obstruction stage cannot be omitted]
\label{C:baseline-free-obstruction-example}
There are compatible crossed-module data with $C=Q=\mathbb Z/2$
for which the obstruction (C.4.3) is nonzero; hence no central-extension
parameter gives a realization.  Explicitly take
\[
 N=L=\mathbb Z/4,\quad d(n)=2n,\quad
 \theta_l(n)=(-1)^l n,\quad Z=C=\{0,2\}\subset N,
 \quad M=Q=\mathbb Z/2,
\]
with $e=0$ and the evident quotient maps.  The nonzero normalized value
of the obstruction is $o(1,1,1)=2\in C$.
\end{proposition}
\begin{proof}
The action fixes $C$, satisfies equivariance of $d$, and is trivial on
$d(N)$, as required because $N$ is abelian.  Here $H=N$, $K=L$.
Take $s(0)=0,s(1)=1$ and $f(1,1)=1$; all normalized remaining values
are zero.  Formula (C.4.3), in additive notation, gives
$-1-1=2\pmod4$.  A normalized $C$-valued $2$-cochain has only the
value $t(1,1)$, and its coboundary at $(1,1,1)$ is
$t(1,1)-t(1,1)=0$.  Thus (C.4.4) has no solution.  The supplied
finite certificate enumerates both cochains, verifies the crossed-module
identities and cocycle identity, and confirms this obstruction.
It also checks all $512$ normalized $C$-valued candidate $2$-cochains
for $Q=(\mathbb Z/2)^2$, $H=C=Z=N=\mathbb Z/2$ and zero boundary;
exactly $16$ are cocycles, and exactly $8$ yield centre and derived group
both equal to $C$.  A separate order-$54$ example has nonabelian $N$,
so the construction has not assumed the derived group to be abelian.
\end{proof}

\subsection{C.4(i): noncommutative unipotent formal kernels}

We specify the category, since the adjective ``smooth'' alone gave the
perfect-field counterexample in Proposition~\ref{C:formal_Frobenius_counterexample}.
A connected formal group of finite type over $k$ is represented by a
complete noetherian local $k$-algebra with residue field $k$; its law need
not be commutative.  Exactness below means exactness as flat sheaves,
with the indicated formal quotients.  We use the usual finite-filtration
meaning of a unipotent formal group: it has a finite closed normal
series with factors isomorphic to closed formal subgroups of additive
formal groups.  The proof in fact needs only that the kernel be
solvable and that some single power $p^a$ kill every group-valued point.
Nonsmooth kernels are included.  Neither a
nilpotent tangent Lie algebra nor pointwise torsion without a uniform
bound is substituted for this hypothesis.

\begin{theorem}[Proved: the noncommutative unipotent alternative]
\label{C:formal-noncommutative-unipotent}
Let $k$ be a perfect field of characteristic $p>0$, and let
\[
 1\longrightarrow U\longrightarrow G\xrightarrow{q}P
 \longrightarrow1
\]
be an extension of connected formal groups of finite type, with $P$ a
connected $p$-divisible group.  Suppose that $U$ is solvable and has
bounded $p$-power exponent; in particular, $U$ may be an arbitrary
noncommutative unipotent formal group in the finite-filtration sense
just specified.  There are an integer $b\geq0$ and a homomorphism
$s:P\to G$ such that $qs=[p^b]_P$.  Equivalently, pullback by this
isogeny of the quotient is a split extension.

Over an algebraically closed field of characteristic $p>0$ the same
conclusion holds, with exactly the same kernel hypotheses.  Over a
perfect field no algebraic-closure or Galois-splitting hypothesis is
needed for this unipotent assertion.  It is a semidirect-product
splitting of the pullback; the resulting action on $U$ need not be
trivial.  All statements include $p=2$.
\end{theorem}

We prove the theorem through three lemmas.  The structural facts about
formal groups used in the reductions are stated for noncommutative
groups in Demazure, II, \S10, pp.~41--43
\cite{C4-Demazure1972}: over a perfect field a reduced connected formal
group of finite type is smooth, and quotients by closed subgroups have
finite-type formal coordinate algebras.  Quotient existence and the
effective torsor identity are supplied by SGA~3, VIIB, 2.4; 2.4.1
identifies a homomorphism's quotient by its kernel with its closed
image \cite{C4-SGA3VIIB}.  Over a field the topological flatness
hypothesis on every subgroup in that result is automatic.  These are
the exact formal statements needed here, rather than the analogous
assertions for algebraic groups.  Smoothness of a high Frobenius image
will also follow directly from the nilradical calculation below.

\begin{lemma}[Proved: a vector kernel]
\label{C:formal-vector-splitting}
Under the field and quotient hypotheses of the theorem, every extension
with kernel $A\simeq\widehat{\mathbb G}_a^d$ admits a lift of $[p]_P$.
If $P$ is multiplicative, the extension itself splits.
\end{lemma}
\begin{proof}
First suppose that $P=T$ is multiplicative.  A finite-dimensional
representation of $T$ is semisimple.  To see this over $\bar k$, write
$T_{\bar k}\simeq\widehat{\mathbb G}_m^r$.  Restriction to
$T[p^n]$ decomposes a finite-dimensional vector space into its
$(\mathbb Z/p^n)^r$ weight spaces.  These decompositions refine with
$n$, and hence stabilize as decompositions into subspaces.  Each
stable subspace carries one compatible system of weights, in
$\mathbb Z_p^r$, and $T$ acts on it by the resulting scalar character.
The subgroup schemes $T[p^n]$ exhaust the formal group, so equality on
all these restrictions is equality of the representations.  A
complement can therefore be chosen in each weight space.  Semisimplicity
over $k$ follows as well: equivariance of a proposed linear splitting
is a system of linear equations in finitely many coefficients over
$k$.  Existence over $\bar k$ implies consistency over $k$, even if the
equations were initially written as an infinite coefficient system.

The conjugation action of $T$ on $A$ can be linearized by an
\emph{additive} change of formal coordinates.  Indeed, primitive
functions on $A$ are the series
\[
  \sum_{i=1}^d\sum_{j\geq0} a_{ij}X_i^{p^j}.
\]
Filter this space by the terms with $j\geq n$.  Each quotient is a
finite-dimensional $T$-representation, and the quotient for $n=1$ is
the cotangent space.  Semisimplicity lets us lift its identity map
successively and compatibly through all these quotients.  The limit
gives $d$ primitive functions with independent linear terms, spanning
a $T$-stable vector space.  They are additive formal coordinates, and
in these coordinates the action is the linear action on a vector
space $V$.

Since the kernel is smooth, $G\to T$ is formally smooth.  Starting at
the identity, lift a section successively over the Artin quotients of
the power-series ring of $T$.  The limit is a normalized formal-scheme
section.  Its defect is a normalized formal $2$-cocycle
$f:T\times T\to\widehat V$ for the linear action.  Put $T_n=T[p^n]$.
A finite group scheme of multiplicative type has a normalized
translation-invariant integral on its coordinate ring: after splitting
it is the projection onto the coefficient of the identity character,
and this projection descends.  In particular, its ordinary regular
cochain complex with coefficients in $V$ has zero cohomology in degrees
$1$ and $2$.  Here one can verify the needed degrees directly.  If
$z(gh)=z(g)+g z(h)$, integration in $h$ gives
$z(g)=a-ga$, where $a=\int z(h)$.  If
\[
 f(g,h)+f(gh,l)=g f(h,l)+f(g,hl),
\]
integration in $l$ gives $f=\delta b$, where
$b(g)=\int f(g,l)$.  Normalized cocycles give $b(1)=0$.

Choose $b_n:T_n\to V$ with $\delta b_n=f|_{T_n^2}$.
They can be chosen compatibly.  Given $b_n$ and any choice
$b'_{n+1}$, the difference
$b'_{n+1}|_{T_n}-b_n$ is a $1$-cocycle, hence is $\delta a_n$ for
some $a_n\in V(k)$.  Replace $b'_{n+1}$ by
$b'_{n+1}-\delta a_n$.  The adjusted cochains restrict to $b_n$.
Although $a_n$ is an algebraic vector-group point, its coboundary
vanishes at the identity and is a legitimate formal cochain.  Thus
this adjustment does not assume that $\widehat V(k)$ contains
$a_n$.  Since the ideals defining $T_n$ are cofinal in the formal
topology, the compatible $b_n$ give a formal $b:T\to\widehat V$.
The identity $\delta b=f$ holds on every finite level and therefore
formally.  Correcting the section by $b$ proves the multiplicative
case, including the compatibility which a finite-level splitting
theorem alone would not supply.

For general $P$, the split connected--etale decomposition over a
perfect field, applied to its Cartier dual, gives
\[
 P=T\times B,
\]
where $T$ is multiplicative and both $B$ and $B^D$ are connected.
This decomposition is integral: Demazure, II, \S7, pp.~34--35 and
\S11, pp.~45--47, give the decomposition and duality used here
\cite{C4-Demazure1972}.  In particular $B$ has no nonzero character
to $\widehat{\mathbb G}_m$, even after extending $k$ to $\bar k$:
duality would give a map from an etale $p$-divisible group to $B^D$.

Every finite-dimensional representation of $B$ is trivial.  For a
proof, take the schematic closure $H$ of its image in a general linear
group after extending to $\bar k$.  Its coordinate algebra embeds in
the reduced formal coordinate algebra of $B$, so $H$ is reduced;
it is connected and commutative.  Every character of $H$ restricts
to zero on $B$, and schematic density makes it zero on $H$.
Simultaneous triangularization of the commuting matrices of $H$
therefore has diagonal entries all equal to $1$.  Thus $H$ is
unipotent and has bounded $p$-power exponent.  Surjectivity of every
$[p^m]_B$ then forces the representation to be trivial.  The conclusion
descends to $k$.

Apply this to the action on every finite coordinate jet of $A$.
The action of $B$ on $A$ is trivial.  Let $G_B$ be its inverse image
in $G$.  The commutator of this central extension descends to a
biadditive map $B\times B\to A$.  It is zero because $pA=0$ and
$[p]_B$ is an epimorphism.  Hence $G_B$ is commutative.  Multiplication
by $p$ on $G_B$ kills $A$ and factors through a homomorphism
$j:B\to G_B$ with $qj=[p]_B$.  This lift is unique, since
$\operatorname{Hom}(B,A)=0$ by the same divisibility argument.
It centralizes all of $G$: conjugating $j$ by any section of $G$
gives another lift, and the difference is a homomorphism to $A$.
This uniqueness argument is valid after every base change, since
$[p]_B$ remains an epimorphism and $pA=0$.

The inverse image $G_T$ splits by the multiplicative case.  Let
$i:T\to G_T$ be a splitting.  As $j(B)$ centralizes $i(T)$, the map
\[
 (t,b)\longmapsto i([p]t)j(b)
\]
is a homomorphism lifting $[p]_{T\times B}$.  This proves the lemma.
\end{proof}

\begin{lemma}[Proved: a characteristic vector filtration]
\label{C:formal-unipotent-filtration}
Let $U$ be a smooth connected solvable formal group of finite type over
a perfect field of characteristic $p$, of bounded $p$-power exponent.
It admits a finite characteristic closed normal filtration with factors
isomorphic to $\widehat{\mathbb G}_a^{d_i}$.  Here characteristic
means preserved by formal automorphism actions after base change, so
the filtration is preserved by conjugation in any formal extension.
\end{lemma}
\begin{proof}
Use the closed derived series.  Each term is smooth: it is the closed
subgroup generated by commutator images and their finite products,
whose sources are products of smooth formal groups.  The vanishing
ideal of this closed image is an intersection of radical ideals, hence
is radical.  A reduced formal group over a perfect field is smooth by
the theorem quoted above.  For invariance under an action with a
possibly nonreduced parameter formal scheme, use coefficient
equations.  If a function vanishes on every product of commutators,
its pullback under universal conjugation vanishes on every such
product with parameters.  Expanding in a topological $k$-basis of the
parameter algebra shows that every coefficient belongs to the same
closed vanishing ideal.  Thus universal conjugation preserves that
ideal, and hence the subgroup.  This does not assume that invariance
under field-valued automorphisms implies invariance under formal
actions.  The derived series terminates by solvability.

The resulting commutative factors are smooth and have bounded
$p$-power exponent.  Refine each by the images of its multiplication
maps $[p^j]$.  These images are reduced, hence smooth, and the
successive quotients are smooth and killed by $p$.  All constructions
are characteristic.  A smooth commutative formal group killed by $p$
is a formal vector group over a perfect field.  Here is the exact
Cartier-module deduction.  Under Zink's $V$-reduced equivalence,
Theorem~4.23, pp.~73--74 \cite{C4-Zink}, $V$ is injective, the module
is $V$-adically complete, and the quotient by $V$ is its finite-dimensional
tangent space.  The relations $VF=FV=p$ (loc. cit., Theorem~4.17,
pp.~71--72) and $pM=0$ give $F=0$.  Choose lifts of a basis of
$M/VM$.  Successive expansion modulo $V$, using perfection to undo
the scalar Frobenius, writes each element uniquely as a convergent
sum of $V^j$ times these basis lifts with scalar coefficients.
With $F=0$ the scalar Cartier operators are additive, so this is
precisely the direct sum of $d$ copies of the Cartier module of
$\widehat{\mathbb G}_a$.  The equivalence identifies the group
with $\widehat{\mathbb G}_a^d$.  This proves the claimed filtration.
\end{proof}

\begin{lemma}[Proved: finite kernels]
\label{C:formal-finite-kernel}
Let $1\to K\to E\to P\to1$ be an extension of connected formal
groups of finite type over a perfect field of characteristic $p>0$,
where $K$ is finite and $P$ is connected and $p$-divisible.
There is a $p$-power isogeny of $P$ whose pullback splits.
\end{lemma}
\begin{proof}
The reduction $E_{\mathrm{red}}$ is a formal subgroup: over a perfect
field the completed tensor product of reduced complete local
$k$-algebras is reduced, so the group law preserves the nilradical.
It is smooth by Demazure's theorem.  The map
$E_{\mathrm{red}}\to P$ is finite and schematically dominant.
Indeed, $\mathcal O(P)\to\mathcal O(E)$ is finite faithfully flat;
its composite with the reduction stays injective because
$\mathcal O(P)$ is reduced.  The two smooth formal groups have the
same dimension.  This finite map is flat: its source is
Cohen--Macaulay and a regular system of parameters in the regular
target is a system of parameters in the finite source, hence a regular
sequence.  The Koszul resolution then gives
$\operatorname{Tor}_1^{\mathcal O(P)}(k,\mathcal O(E_{\mathrm{red}}))=0$;
lifting a residue-field basis and applying Nakayama to the relation
module makes this finite module free.  Thus the map is an isogeny of
formal groups.

The commutator on $E_{\mathrm{red}}$ takes values in $K$.  A map from
a connected smooth formal scheme to a finite connected scheme factors
through the latter's reduced identity point, so this commutator is
zero.  Therefore $E_{\mathrm{red}}$ is commutative.  The finite kernel
of $E_{\mathrm{red}}\to P$ is a commutative finite connected group and
is killed by some $p^c$ (Demazure, II, \S9, p.~40).
Multiplication by $p^c$ on $E_{\mathrm{red}}$ descends along the isogeny
to a map $P\to E_{\mathrm{red}}$ whose composite to $P$ is
$[p^c]_P$.  Inclusion into $E$ gives the required lift.
\end{proof}

\begin{proof}[Proof of Theorem~\ref{C:formal-noncommutative-unipotent}]
Suppose first that $U$ is smooth, and choose the characteristic
filtration of Lemma~\ref{C:formal-unipotent-filtration}.  Induct on
its length $r$.  For $r=0$ there is nothing to prove.  Write its first
nonzero term as $U_1\simeq\widehat{\mathbb G}_a^d$; it is normal
in $G$.  By induction the extension $G/U_1\to P$ admits a lift of
$[p^{r-1}]_P$.  Pulling $G\to G/U_1$ back along that lift gives an
extension of $P$ by $U_1$.  Lemma~\ref{C:formal-vector-splitting}
gives a lift of $[p]_P$ to this pullback.  Its composite to $G$ lifts
$[p^r]_P$.  No commutativity of $U$ was used.

For an arbitrary $U$, choose $n$ such that $p^n$-th powers kill its
nilradical.  The kernel of the corresponding Frobenius map on its
coordinate algebra is then exactly that nilradical.  The image of
$F_U^n$ therefore has reduced coordinate algebra (with the appropriate
Frobenius twist) and is smooth by the noncommutative reduced-group
theorem.  Frobenius is finite on a complete finite-type local algebra
over a perfect field, so this is a closed formal image; equivalently,
use the quotient--image assertion of SGA~3, VIIB, 2.4.1.
Its kernel $K$ is finite: the $p^n$-powers of finitely many generators
of the maximal ideal define an ideal of finite colength.
It is preserved by conjugation in $G$, since Frobenius commutes with
conjugation.  The quotient $U/K$ is smooth, solvable and of bounded
$p$-power exponent.  Apply the smooth case to $G/K\to P$ and pull
$G\to G/K$ back along the resulting lift.  This is an extension
of $P$ by the finite kernel $K$.  Lemma~\ref{C:formal-finite-kernel}
supplies another $p$-power lift; composing the two lifts gives $s$.

A finite unipotent filtration gives the hypotheses on $U$: additive
factors are abelian and killed by $p$, so a series of length $a$
gives solvability and the uniform identity $u^{p^a}=1$.
Thus the theorem decides the unipotent alternative of scan~36(I)
in this finite-type formal category.  The multiplicative Galois
counterexample has no such bounded kernel and therefore places no
restriction on this conclusion over perfect fields.
\end{proof}

